\documentclass[10pt,a4paper]{amsart}
\usepackage[margin=19mm]{geometry}
\usepackage{amsmath,amssymb,mathtools}
\usepackage[scr=rsfs,cal=euler]{mathalfa}
\usepackage{xfrac}
\usepackage[unicode,hidelinks,bookmarks=false]{hyperref}
\newcommand{\ie}{i.e.\ }
\newcommand{\cf}{cf.\ }
\newcommand{\resp}{respectively\ }
\newcommand{\Subsection}{\subsection}
\providecommand{\nobreakdash}{\nobreak\hbox{-}}
\setlength{\emergencystretch}{2em}
\multlinegap0pt
\usepackage{bm}
\usepackage{bbm}
\usepackage{dsfont}
\usepackage{xspace}
\usepackage{cancel}
\usepackage{mathtools}
\usepackage{amsthm}
\makeatletter
\@ifundefined{thm}{\newtheorem{thm}{Thm}}{}
\@ifundefined{defn}{\newtheorem{defn}[thm]{Definition}}{}
\@ifundefined{lem}{\newtheorem{lem}[thm]{Lemma}}{}
\makeatother
\newcommand{\E}{\mathbb{E}}
\renewcommand{\P}{\mathbb{P}}
\title[The Poisson boundary of wreath products]{The Poisson boundary of wreath products}
\author{Joshua Frisch, Eduardo Silva}
\date{Source: arXiv:2310.10160v4 (CC BY 4.0)}
\begin{document}
\maketitle

\noindent\emph{Source and licence.} The Poisson boundary of wreath products. Version arXiv:2310.10160v4 (CC BY 4.0), distributed under the Creative Commons Attribution 4.0 International licence, as recorded in the arXiv licence metadata for this exact version and shown on the abstract page https://arxiv.org/abs/2310.10160v4. Full rights notice: licenses/arxiv-2310.10160-CC-BY-4.0.txt.\par\smallskip

\noindent\textit{Editorial scope.} Counted unit: the Lem whose source label is \texttt{lem: expectation of unstable points is small} in the article (source0.tex lines 675--678, source label \texttt{lem: expectation of unstable points is small}), delivered together with the article's own complete original proof of it (source0.tex lines 679--718, one \texttt{proof} environment of 40 lines). No mathematics is written, repaired or re-derived: statement and proof are the article's own text line for line. Required context retained uncounted: def:neighborhood of coarse traj (lines 575--580). Every definition, display and label of those lines is retained unchanged. Omitted from this package and not redistributed: 1--574, 581--674, 719--1343. Nothing inside a retained excerpt is omitted or altered: every retained body is delivered exactly as the article's lines are written, so no omission receipt of any kind accompanies this package. Citations: the retained excerpts carry no citation command at all, so no bibliography entry of the article is transcribed and no reference is redistributed. Reference adaptation: none was needed and none was made; every reference inside the delivered unit resolves inside it, so no unresolved reference or citation is produced. Printed slips kept literal: no printed word, symbol, hypothesis or number of the counted statement or of its proof was altered. Layout and class: the article's own class is \texttt{amsart} with packages including fontenc, inputenc, lmodern, babel, geometry, amsmath, amssymb, amsfonts, amsthm, mathtools, mathrsfs, graphicx, caption, subcaption; the wrapper is the lane's pinned \texttt{amsart} editorial wrapper at 10pt on A4, which restates the theorem-like environments the retained text needs and adds the article's own macro definitions that the retained text needs (\texttt{\textbackslash E}, \texttt{\textbackslash P}), with extra wrapper packages (bm, bbm, dsfont, xspace, cancel, mathtools). The delivered numbering is the unit's own. Assets: no retained span contains a figure environment, a graphics-inclusion command, a PDF-page inclusion, a table or a TikZ picture; no image file is copied into this package, the package declares no asset dependency and no image file is redistributed with it. Licence: version arXiv:2310.10160v4 of this article is distributed under the Creative Commons Attribution 4.0 International licence, as recorded in the arXiv licence metadata for this exact version and shown on the abstract page https://arxiv.org/abs/2310.10160v4. A full rights notice is delivered at licenses/arxiv-2310.10160-CC-BY-4.0.txt. Characters: every retained excerpt is pure ASCII, so the delivered engine, XeLaTeX with its default Latin Modern text font, reports no missing character.
\begin{defn}\label{def:neighborhood of coarse traj}
	Let $t_0\ge 1$ and $R\subseteq B$ be a finite subset with $e_B\in R$. Define the \emph{$(t_0,R)$-coarse neighborhood} of the trajectory at instant $n$ by
	\begin{equation*}
		\mathcal{N}_n(t_0,R)\coloneqq \bigcup_{j=0}^{\lfloor n/t_0 \rfloor -1} \left\{ X_{jt_0}r_1r_2\cdots r_{t_0}\mid r_k\in R, \text{ for }k=1,2,\ldots,t_0 \right \}.
	\end{equation*}
\end{defn}

\begin{lem}\label{lem: expectation of unstable points is small} Let $t_0\ge 1$ and consider finite subsets $R\subseteq B$ and $L\subseteq A$ with $e_B\in R$ and $e_A\in L$. For any $\varepsilon>0$, there exists a constant $K_2>0$ such that for every $n\ge 1$, 
	\begin{equation*}\E\left(|\mathcal{U}_s| \mid \mathcal{P}_n^{t_0}\right)<\varepsilon n+K_2, \text{ where we denote }s=\lfloor n/t_0\rfloor t_0.
	\end{equation*}
\end{lem}

\begin{proof}
	Consider $t_0,R$ and $L$ as above. For $0\le j\le \lfloor n/t_0\rfloor-1$, let us denote $F_j\coloneqq X_{jt_0}R^{t_0}$ and recall that $\mathcal{N}_n(t_0,R)=\bigcup_{j=0}^{\lfloor n/t_0\rfloor-1}F_j$ (Definition \ref{def:neighborhood of coarse traj}).
	
	For $n_0\ge 1$, let us say that an instant $j$ is \emph{poorly $n_0$-stabilized} if the lamp configuration at $F_j$ is modified at some instant beyond $jt_0+n_0$. That is, if for some $k>jt_0+n_0$ and some $b\in F_j$, we have $\varphi_{k-1}(b)\neq \varphi_{k}(b)$.
	
	We observe the following.\\
	\textbf{Claim 1:} \textit{We have that \( \P(j \text{ is poorly }n_0\text{-stabilized})=\P(0 \text{ is poorly }n_0\text{-stabilized})\) for every $n_0\ge 1$ and $j\ge 0$.}
	
	
	Indeed, we see that
	\begin{align*}
		\P(j \text{ is poorly }n_0\text{-stabilized})&=\P(\varphi_{k-1}(b)\neq \varphi_{k}(b)\text{ for some }k>jt_0+n_0\text{ and }b\in X_{jt_0}R^{t_0})\\
		&=\P(\varphi_{k-1}(b)\neq \varphi_{k}(b)\text{ for some }k>n_0\text{ and }b\in R^{t_0})\\
		&=\P(0 \text{ is poorly }n_0\text{-stabilized}).
	\end{align*}
	In the second equality above we used the Markov property and the group invariance of the random walk to change the time instants beyond $jt_0+n_0$ by those beyond $n_0$.
	
	\textbf{Claim 2:}\textit{ $\lim_{n_0\to \infty}\P(0 \text{ is poorly }n_0\text{-stabilized})=0$.}\\
	Indeed, this follows from the hypothesis that lamp configurations stabilize almost surely, which implies that $\P(\text{the lamp configuration in }R^{t_0} \text{ is modified after time $n_0$})\xrightarrow[n_0\to \infty]{}0.$
	
	Let $\varepsilon>0$. Using Claims 1 and 2 we can find $n_0\ge 1$ such that for every $k\ge n_0$ and $j\ge 0$ we have $
		\P(j \text{ is poorly }k\text{-stabilized})\le \P(j \text{ is poorly }n_0\text{-stabilized})<\varepsilon.$
	
	Thanks to the linearity of expectation, we have
	\begin{align*}
		\E\left(|\mathcal{U}_s| \mid \mathcal{P}_n^{t_0}\right)&= 
		\E\left(\sum_{j=0}^{\lfloor n/t_0\rfloor-1}\mathds{1}_{\left\{\varphi_s|_{F_j}\neq \varphi_{\infty}|_{F_j}\right \}} \mid \mathcal{P}_n^{t_0}\right)\\
		&=
		\sum_{j=0}^{\lfloor n/t_0\rfloor-1}\P\left(\varphi_s|_{F_j}\neq \varphi_{\infty}|_{F_j}\mid \mathcal{P}_n^{t_0}\right)=\sum_{j=0}^{\lfloor n/t_0\rfloor-1} \P(j\text{ is poorly }s\text{-stabilized}).
	\end{align*}
	
	Denote $J=\left \lfloor \lfloor n/t_0 \rfloor- n_0/t_0-1  \right \rfloor$, which is the largest integer that satisfies the inequality $Jt_0+t_0\le \lfloor n/t_0 \rfloor t_0 - n_0$. In other words, $Jt_0$ is the last instant whose associated interval $I_{Jt_0}$ still has $n_0$ remaining units of time before reaching time $s$. Then if $j\in \{0,1,\ldots, J\}$, it holds that $  \P(j\text{ is poorly }s\text{-stabilized})\le \P(0\text{ is poorly }n_0\text{-stabilized}).$ From this, we see that
	\begin{align*}
		\E\left(|\mathcal{U}_s| \mid \mathcal{P}_n^{t_0}\right)&\le\sum_{j=0}^{J} \P(j\text{ is poorly }s\text{-stabilized})+\sum_{j=J+1}^{\lfloor n/t_0\rfloor-1}\P(j\text{ is poorly }s\text{-stabilized})\\
		&\le (J+1) \P(0\text{ is poorly }n_0\text{-stabilized})+ \lfloor n/t_0\rfloor-1 -J \\
		&< (J+1)\varepsilon +\lfloor n/t_0\rfloor-1 -J.
	\end{align*}
	Our choice of $J$ guarantees that $J+1\le (n-n_0)/t_0\le n$ and $\lfloor n/t_0\rfloor-1 -J\le n_0/t_0+1$, so that $
		\E\left(|\mathcal{U}_s| \mid \mathcal{P}_n^{t_0}\right)\le \varepsilon n + K_2,$ where $K_2=n_0/t_0+1$ is a constant that only depends on $\mu$, $t_0$, $R$ and $\varepsilon$.
\end{proof}

\end{document}
